\exercisesection{Exercises for § 1}

All the root systems considered below are relative to real vector spaces. We denote by \((x|y)\) a scalar product invariant under the Weyl group (cf.~no. 3).

\begin{enumerate}
\def\labelenumi{\arabic{enumi})}
\item
  Let \(\mathbb{R}\) be a root system, \(\mathbb{R} = \mathbb{R}_1 \cup \mathbb{R}_2\) a partition of \(\mathbb{R}\) . Assume that, if \(x, y\) are two elements of \(\mathbb{R}_i\) and if \(x + y\) (resp. \(x - y\) ) is a root, then \(x + y \in \mathbb{R}_i\) (resp. \(x - y \in \mathbb{R}_i\) ), where \(i = 1, 2\) . Then, \(\mathbb{R}\) is the direct sum of \(\mathbb{R}_1\) and \(\mathbb{R}_2\) . (By using the Cor. of Th. 1, no. 3, show that if \(x \in \mathbb{R}_1\) and \(y \in \mathbb{R}_2\) then \((x|y) = 0\) .)
\item
  Let R be a root system, \(\alpha\) and \(\beta\) two roots. If t is a scalar such that \(\beta + t\alpha \in R\) , then \(2t \in Z\) . (For \(n(\beta + t\alpha, \alpha) = n(\beta, \alpha) + 2t\) .) If \(\alpha\) is indivisible, then \(t \in Z\) . (If not, show, by using Prop. 9, that there exists a root \(\gamma\) orthogonal to \(\alpha\) such that \(\gamma + \frac{1}{2}\alpha \in R\) : then \((\gamma + \frac{1}{2}\alpha|\alpha) > 0\) , so \(\frac{1}{2}\alpha \in R\) .)
\item
  Let R be an irreducible non-reduced root system in V. There exists a non-degenerate symmetric bilinear form \(\Phi\) on V such that, if we identify V with \(V^{*}\) using \(\Phi\) , we have \(R = R^{-}\) . (Use Prop. 13.)
\item
  Let \(\Gamma\) be a free \(\mathbf{Z}\) -module of finite rank \(l\) , \(\Gamma^{*}\) the dual \(\mathbf{Z}\) -module, I a finite set of indices, \((x_{i}, x_{i}^{*})_{i \in \mathbf{I}}\) a family of elements of \(\Gamma \times \Gamma^{*}\) such that \(\langle x_{i}, x_{i}^{*} \rangle = 2\) for all \(i \in \mathbf{I}\) , \(s_{i} = s_{x_{i}, x_{i}^{*}}\) . Let \(\mathbf{V}\) be the real vector space \(\Gamma \otimes_{\mathbf{Z}} \mathbf{R}\) , whose dual \(\mathbf{V}^{*}\) can be identified with \(\Gamma^{*} \otimes_{\mathbf{Z}} \mathbf{R}\) . We also denote by \(s_{i}\) the reflection \(s_{i} \otimes 1\) in \(\mathbf{V}\) . Let \(\mathbf{R}\) (resp. \(\mathbf{R}'\) ) be the set of \(x_{i}\) (resp. \(x_{i}^{*}\) ). Let \(\mathbf{E}\) (resp. \(\mathbf{E}'\) ) be the subspace of \(\mathbf{V}\) (resp. \(\mathbf{V}^{*}\) ) generated by \(\mathbf{R}\) (resp. \(\mathbf{R}'\) ). Assume that \(s_{i}(\mathbf{R}) = \mathbf{R}\) for all \(i\) . Then, \(\mathbf{R}\) is a root system in \(\mathbf{E}\) . \(\mathbf{E}'\) is canonically identified with the dual of \(\mathbf{E}\) and \(x_{i}^{*}\) with \(x_{i}^{-}\) for all \(i\) . If \(\mathbf{F}\) (resp. \(\mathbf{F}'\)) is the subspace of \(\mathbf{V}\) (resp. \(\mathbf{V}^{*}\) ) consisting of the points invariant under all the \(s_{i}\) (resp. \(^t s_{i}\) ), then \(\Gamma \cap \mathbf{F}\) (resp. \(\Gamma^{*} \cap \mathbf{F}'\) ) generates \(\mathbf{F}\) (resp. \(\mathbf{F}'\) ). If we put \(\Gamma_{1} = (\Gamma \cap \mathbf{E}) \oplus (\Gamma \cap \mathbf{F})\) (resp. \(\Gamma_{1}^{*} = (\Gamma^{*} \cap \mathbf{E}') \oplus (\Gamma^{*} \cap \mathbf{F}')\) ), then \(\Gamma / \Gamma_{1}\) and \(\Gamma^{*} / \Gamma_{1}^{*}\) are isomorphic finite groups.
\item
  Let \(R\) be an irreducible root system. For any \(\beta \in R\) ,
\end{enumerate}

\[
16 \gamma (\mathrm{R}) = (\sum_ {\alpha \in \mathrm{R}} n (\alpha , \beta) ^ {2}) (\sum_ {\alpha \in \mathrm{R}} n (\beta , \alpha) ^ {2})
\]

and consequently \(\gamma (\mathbb{R}) = \gamma (\lambda R)\) for any non-zero scalar \(\lambda\) . (Use formulas (17) and (19) of no. 12.)

\begin{enumerate}
\def\labelenumi{\arabic{enumi})}
\setcounter{enumi}{5}
\tightlist
\item
  \begin{enumerate}
  \def\labelenumii{\alph{enumii})}
  \tightlist
  \item
    Let A be an abelian group of finite type, and T a finite subset of A not containing 0. There exists a subgroup H of A of finite index such that \(H \cap T = \varnothing\) . (Let \(t \in T\) . By using the structure of abelian groups of finite type, construct a subgroup of A of finite index that does not contain t. Now proceed by induction on \(\text{Card}(T)\) .)
  \end{enumerate}
\end{enumerate}

\begin{enumerate}
\def\labelenumi{\alph{enumi})}
\setcounter{enumi}{1}
\tightlist
\item
  Let R be a root system and P a closed symmetric subset of R. There exists a subgroup H of Q(R) of finite index such that \(P = H \cap R\) . (Pass to the quotient by the subgroup of Q(R) generated by P, and use a) and Prop. 23.)
\end{enumerate}

\begin{enumerate}
\def\labelenumi{\arabic{enumi})}
\setcounter{enumi}{6}
\item
  The connection index of a root system is equal to the determinant of its Cartan matrix. (Use formula (14) of no. 10. Show on the other hand that, in a real vector space with a scalar product, if \((\varepsilon_1, \ldots, \varepsilon_n), (\varepsilon_1', \ldots, \varepsilon_n')\) are two bases such that \((\varepsilon_i | \varepsilon_j') = \delta_{ij}\) , then \(\det_{(\varepsilon_1, \ldots, \varepsilon_n)} (\varepsilon_1', \ldots, \varepsilon_n') > 0\) .)
\item
  Let R be a reduced root system, C a chamber, \(2\sigma\) the sum of the positive elements of \(R^{-}\) , \(P'(R)\) the set of \(x \in P(R)\) such that \(\langle x, \sigma \rangle \in Z\) . Then, \(Q(R) \subseteq P'(R) \subseteq P(R)\) , and \(P(R)/P'(R)\) is of order 1 or 2. Let \((\beta_{1}, \ldots, \beta_{l})\) be the basis of \(R^{-}\) corresponding to C, and put \(2\sigma = n_{1}\beta_{1} + \cdots + n_{l}\beta_{l}\) where the \(n_{i}\) are integers \textgreater0. Then, \(P(R) = P'(R)\) if and only if all the \(n_{i}\) are even.
\item
  Let \(\mathbb{R}\) be a root system, \(\mathbb{C}\) a chamber. Put \(\mathrm{B}(\mathbb{C}) = \{\alpha_1, \ldots, \alpha_l\}\) ,
\end{enumerate}

\[
\mathrm{R} _ {+} (\mathrm{C}) = \left\{\alpha_ {1}, \dots , \alpha_ {l}, \alpha_ {l + 1}, \dots , \alpha_ {s} \right\}
\]

(the \(\alpha_{i}\) being pairwise distinct), and \(\alpha = \alpha_{1} + \dots +\alpha_{s}\) .

\begin{enumerate}
\def\labelenumi{\alph{enumi})}
\tightlist
\item
  For \(i = 1, 2, \ldots, s\) , put \(\varepsilon_{i} = \pm 1\) . Let \(\alpha^{*} = \varepsilon_{1}\alpha_{1} + \cdots + \varepsilon_{s}\alpha_{s}\) . If \((\alpha^{*}|\alpha_{i}) > 0\) for \(i = 1, \ldots, l\) , then \(\alpha^{*} = \alpha\) and \(\varepsilon_{i} = 1\) for all i. (Let \(\gamma\) (resp. \(\delta\) ) be the sum of the \(\alpha_{i}\) for which \(\varepsilon_{i} = 1\) (resp. -1). Then, \((\gamma|\alpha_{i}) - (\delta|\alpha_{i}) > 0\) , \((\gamma|\alpha_{i}) + (\delta|\alpha_{i}) = (\alpha_{i}|\alpha_{i})\) , so
\end{enumerate}

\(2(\gamma|\alpha_{i})>(\alpha_{i}|\alpha_{i});\) since \(2(\gamma|\alpha_{i}) / (\alpha_{i}|\alpha_{i})\in \mathbf{Z},(\gamma |\alpha_{i})\geqslant (\alpha_{i}|\alpha_{i}) = (\alpha |\alpha_{i}),\) so \(\gamma -\alpha \in \overline{\mathrm{C}}. \gamma \geqslant \alpha ,\) \(\gamma = \alpha\) and \(\delta = 0.)\)

\begin{enumerate}
\def\labelenumi{\alph{enumi})}
\setcounter{enumi}{1}
\tightlist
\item
  For \(i = 1,2,\ldots,s\) , let \(\varepsilon_{i} = \pm 1\) . The set of \(\varepsilon_{i}\alpha_{i}\) is the set of roots \(>0\) relative to a chamber if and only if \(\alpha^{*} = \sum_{i}\varepsilon_{i}\alpha_{i}\) belongs to a chamber. (To show that the condition is sufficient, let \(\mu_{i} = \pm 1\) be such that \((\alpha^{*}|\mu_{i}\alpha_{i}) > 0\) . There exists an element \(w\in W(R)\) that transforms the set of \(\mu_{i}\alpha_{i}\) into the set of \(\alpha_{i}\) , and hence \(\alpha^{*}\) into \(\alpha^{**} = \sum_{i}\varepsilon_{i}\mu_{i}\alpha_{\sigma(i)}\) , where \(\sigma \in \mathfrak{S}_s\) . Applying \(a\) ), show that \(\alpha^{**} = \alpha\) , so \(\mu_{i}\varepsilon_{i} = 1\) .)
\end{enumerate}

\begin{enumerate}
\def\labelenumi{\arabic{enumi})}
\setcounter{enumi}{9}
\item
  Let \(\alpha\) be the highest root of an irreducible reduced root system R relative to some basis. Then, \(\alpha^{-}\) is the highest root of \(R^{-}\) if and only if all the roots are of the same length.
\item
  With the notation of Prop. 33 of no. 11. show that \(\theta_{i} + \theta_{j}\) is not a root for any pair \((i,j)\) .
\item
  Let R be a root system in V of rank \(\geqslant 3\) , \((\alpha_{1},\ldots,\alpha_{l})\) a basis of R. \(V'\) (resp. \(V''\) ) the subspace of V generated by the \(\alpha_{i}\) for \(i \geqslant 2\) (resp. \(i \geqslant 3\) ). \(R' = R \cap V'\) , \(R'' = R \cap V''\) . Let d (resp. \(d', d''\) ) be the determinant of the Cartan matrix of R (resp. \(R', R''\) ). Assume that \(\alpha_{1}\) is orthogonal to all the \(\alpha_{i}\) except \(\alpha_{2}\) , and that \(\| \alpha_{1} \| = \| \alpha_{2} \|\) . Show that \(d = 2d' - d''\) .
\item
  Let R be a reduced root system. \((\alpha_{1},\ldots,\alpha_{l})\) a basis of R and \(\alpha=c_{1}\alpha_{1}+\cdots+c_{l}\alpha_{l}\) a root. Then, \(c_{i}(\alpha_{i}|\alpha_{i})/(\alpha|\alpha)\in\mathbf{Z}\) for all i. (Consider the inverse system.)
\item
  Let R be an irreducible root system, \(\lambda\) the greatest root length. S the set of subsets of R consisting of pairwise orthogonal roots of length \(\lambda\) . Then, any two maximal elements of S are transformed into each other by W(R). (Use Prop. 11. and Prop. 1 of Chap. V. §3. no. 3.)
\item
  Let \(R\) be a root system of rank \(l\) .
\end{enumerate}

\begin{enumerate}
\def\labelenumi{\alph{enumi})}
\item
  \(-1 \in \mathrm{W}(\mathrm{R})\) if and only if R contains l roots that are pairwise strongly orthogonal. (To show that the condition is necessary, argue by induction on l using Prop. 1 of Chap. V, §3, no. 3.)
\item
  If \(w \in \mathrm{W}(\mathrm{R})\) is of order 2, there exists a set S of pairwise strongly orthogonal roots such that w is the product of the \(s_{\alpha} (\alpha \in \mathrm{S})\) .
\end{enumerate}

\begin{enumerate}
\def\labelenumi{\arabic{enumi})}
\setcounter{enumi}{15}
\tightlist
\item
  Let \(\mathbb{R}\) be a root system. B a basis of \(\mathbb{R}\) , and \(\mathbb{R}_+\) the set of positive roots relative to this basis. For \(w \in W(\mathbb{R})\) , put \(F_w = R_+ \cap w(-R_+)\) . Show that the map \(w \mapsto F_w\) is a bijection from \(W(\mathbb{R})\) to the set of subsets \(F\) of \(R_+\) such that \(F\) and \(R_+ - F\) are closed. (Apply Cor. 1 of Prop. 20 to the subset \(P = (R_+ - F) \cup (-F)\) .)
\end{enumerate}

¶ 17) Let R be a reduced root system and B a basis of R. For any subset J of B, let W \(_{J}\) be the subgroup of W(R) generated by the reflections s \(_{\alpha}\) for \(\alpha \in J\) , and let w \(_{J}\) be the longest element of W \(_{J}\) . Let \(\Phi\) be the set of w \(_{J}\) .

\begin{enumerate}
\def\labelenumi{\alph{enumi})}
\tightlist
\item
  Show that an element \(w \in \mathrm{W}(\mathbb{R})\) belongs to \(\Phi\) if and only if
\end{enumerate}

\[
\mathrm{W} (\mathrm{B}) \subseteq \mathrm{R} _ {+} (\mathrm{B}) \cup (- \mathrm{B}).
\]

(To show that the condition is sufficient. denote by \(\mathrm{J}_1\) the set of \(\alpha \in \mathrm{B}\) such that \(w(\alpha) \in -\mathrm{B}\) and put \(\mathrm{J} = -w(\mathrm{J}_1)\) . If \(\alpha \in \mathrm{J}_1\) , \(w(\alpha) \in -\mathrm{J}\) and \(w_{\mathrm{J}}w(\alpha) \in \mathrm{B}\) ;

if \(\alpha \in \mathrm{B} - \mathrm{J}_1\) , \(w_{\mathrm{J}}w(\alpha) \in \mathrm{R}_{+}\) : otherwise \(w(\alpha)\) would be positive and \(w_{\mathrm{J}}w(\alpha)\) negative, which would imply that \(w(\alpha)\) belongs to the subsystem generated by the \(\beta \in J\) and that \(\alpha\) belongs to the subsystem generated by the \(\beta \in \mathrm{J}_1\) , which is absurd. Deduce that \(w_{\mathrm{J}}w = 1\) .

\begin{enumerate}
\def\labelenumi{\arabic{enumi})}
\setcounter{enumi}{17}
\item
  Let R be a root system and let P be a parabolic subset of R. Show that the complement of P in R is closed.
\item
  Let \(\mathbb{R}\) be a root system, and let \(x\) be a non-zero element of \(\mathrm{Q}(\mathbb{R})\) of minimal length. Show that \(x \in \mathbb{R}\) .
\end{enumerate}

¶ 20) Let \((\alpha_{1},\ldots,\alpha_{l})\) be a basis of an irreducible reduced root system R. Let r and p be two integers, with \(p\geqslant2\) , such that:

\[
\left(\alpha_ {1} \mid \alpha_ {1}\right) = \dots = \left(\alpha_ {r} \mid \alpha_ {r}\right) = p. \left(\alpha_ {r + 1} \mid \alpha_ {r + 1}\right) = \dots = p. \left(\alpha_ {l} \mid \alpha_ {l}\right).
\]

\begin{enumerate}
\def\labelenumi{\alph{enumi})}
\item
  Let \(\alpha = c_{1}\alpha_{1} + \cdots + c_{l}\alpha_{l}\) be a root. Show that \((\alpha|\alpha) = (\alpha_{1}|\alpha_{1})\) (in other words, that \(\alpha\) is a long root) if and only if p divides \(c_{r+1},\ldots,c_{l}\) .
\item
  Let h be the Coxeter number of W(R). Show that the number of longest (resp. shortest) roots of R is equal to hr (resp. \(h(l - r)\) ).
\end{enumerate}

\begin{enumerate}
\def\labelenumi{\arabic{enumi})}
\setcounter{enumi}{20}
\tightlist
\item
  Let R be a root system and B a basis of R. Let \(\alpha, \beta \in \mathrm{B}\) and \(w \in \mathrm{W}(\mathrm{R})\) be such that \(\beta = w(\alpha)\) . Show that there exists a sequence \(\alpha_1, \ldots, \alpha_n\) of elements of R and a sequence \(w_1, \ldots, w_{n-1}\) of elements of \(\mathrm{W}(\mathrm{R})\) such that
\end{enumerate}

\begin{enumerate}
\def\labelenumi{(\roman{enumi})}
\item
  \(\alpha_{1} = \alpha, \quad \alpha_{n} = \beta.\)
\item
  \(w = w_{n - 1}\dots w_1\)
\item
  \(w_{i}(\alpha_{i}) = \alpha_{i + 1}\) for \(1\leqslant i\leqslant n - 1\)
\item
  For all \(i \in [1, n - 1]\) , there exist \(\beta_i \in \mathbf{B}\) such that \(w_i\) belongs to the subgroup of \(\mathrm{W}(\mathbb{R})\) generated by the reflections \(s_{\alpha_i}\) and \(s_{\beta_i}\) .
\end{enumerate}

¶ 22) With the assumptions and notations of Prop. 33 of no. 11, denote by \(\omega_{1},\ldots,\omega_{l}\) the fundamental weights of R.

\begin{enumerate}
\def\labelenumi{\alph{enumi})}
\item
  Show that \(c^{-1}(\omega_i) = \omega_i - \theta_i\) . Deduce that \(1 - c\) maps \(P(R)\) to \(Q(R)\) .
\item
  Let \(f\) be the connection index of \(\mathbb{R}\) , and let \(m_1, \ldots, m_l\) be the exponents of \(W(\mathbb{R})\) . Show that
\end{enumerate}

\[
f = \det (1 - c) = \prod_ {j = 1} ^ {l} \left(1 - \omega ^ {m _ {j}}\right) = 2 ^ {l} \prod_ {j = 1} ^ {l} \sin \left(m _ {j} \pi / h\right)
\]

with \(\omega = e^{2\pi i/h}\) .

\begin{enumerate}
\def\labelenumi{\alph{enumi})}
\setcounter{enumi}{2}
\tightlist
\item
  Let p be a prime number. Write h in the form \(h = p^{a}H\) , with H not divisible by p.~Show that p divides f if and only if H divides one of the \(m_{j}\) .
\end{enumerate}

\begin{enumerate}
\def\labelenumi{\arabic{enumi})}
\setcounter{enumi}{22}
\tightlist
\item
  Let R be a reduced root system, and let X be a subset of P(R). Say that X is saturated if the following condition is satisfied:
\end{enumerate}

\begin{enumerate}
\def\labelenumi{(\Alph{enumi})}
\setcounter{enumi}{18}
\tightlist
\item
  For all \(p \in \mathbf{X}, \alpha \in \mathbb{R}\) and \(i \in \mathbf{Z}\) such that \(i\) is between 0 and \(\langle p, \alpha \rangle\) , \(p - i\alpha \in \mathbf{X}\) .
\end{enumerate}

\begin{enumerate}
\def\labelenumi{\alph{enumi})}
\item
  Show that every saturated subset is stable under the Weyl group W of R.
\item
  Show that, for any subset A of P(R), there exists a smallest saturated subset S(A) of P(R) containing A.
\item
  Let C be a chamber of R, and let \(p \in \mathrm{P}(\mathrm{R}) \cap \overline{\mathrm{C}}\) be a dominant weight. Let \(S(p)\) be the smallest saturated subset of \(\mathrm{P}(\mathrm{R})\) containing p.~Let \(\Sigma(p)\) be the set of elements \(p' \in \mathrm{P}(\mathrm{R})\) such that
\end{enumerate}

\begin{enumerate}
\def\labelenumi{(\roman{enumi})}
\item
  \(p \equiv p' \mod Q(R)\) .
\item
  For all \(w \in \mathbf{W}\) , \(w(p') \leqslant w\) (with respect to the order relation defined by C).
\end{enumerate}

Show that \(\Sigma(p)\) is finite, contains p, and is saturated. Conclude that \(S(p)\) is contained in \(\Sigma(p)\) .

\begin{enumerate}
\def\labelenumi{\alph{enumi})}
\setcounter{enumi}{3}
\tightlist
\item
  Show that, if \(\alpha\) is a longest element of R. \(S(\alpha) = R \cup \{0\}\) .
\end{enumerate}

\begin{enumerate}
\def\labelenumi{\arabic{enumi})}
\setcounter{enumi}{23}
\tightlist
\item
  We retain the notation and assumptions of the preceding exercise.
\end{enumerate}

\begin{enumerate}
\def\labelenumi{\alph{enumi})}
\tightlist
\item
  Let \(p \in \mathrm{P}(\mathbb{R})\) , and let \(\mathbf{W}.p\) be the orbit of \(p\) under \(\mathbf{W}\) . Show that the following conditions are equivalent:
\end{enumerate}

\begin{enumerate}
\def\labelenumi{(\roman{enumi})}
\item
  \(S(p) = W.p.\)
\item
  \(\langle p, \alpha^{\prime} \rangle = 0, 1\) or \(-1\) for all \(\alpha \in \mathbb{R}\) .
\end{enumerate}

(If (ii) is not satisfied, construct an element \(q \in S(p)\) such that \((q|q) < (p|p)\) , and hence \(q \notin W.p.\) )

\begin{enumerate}
\def\labelenumi{\alph{enumi})}
\setcounter{enumi}{1}
\item
  Let X be a non-empty saturated subset of P(R). Show that X contains an element p satisfying conditions (i) and (ii) above. (Take p to be an element of X of minimal length.)
\item
  Assume that \(\mathbb{R}\) is irreducible. Let \(\mathbb{C}\) be a chamber of \(\mathbb{R}\) , let \(\mathbb{B} = \{\alpha_i\}_{i\in I}\) be the corresponding basis, and let
\end{enumerate}

\[
\gamma^ {*} = \sum_ {i} n _ {i} \alpha_ {i}
\]

be the highest root of \(R'\) . Let J be the set of \(i \in I\) such that \(n_{i} = 1\) . Let \(p \neq 0\) be a dominant weight. Show that conditions (i) and (ii) of a) are equivalent to each of the following conditions:

\begin{enumerate}
\def\labelenumi{(\roman{enumi})}
\setcounter{enumi}{2}
\item
  \(\langle p, \gamma^{*} \rangle = 1\) .
\item
  There exists \(i \in \mathrm{J}\) such that \(p\) is equal to the corresponding fundamental weight.
\end{enumerate}

A weight p satisfying these conditions is called minuscule. Show that every non-empty saturated subset of \(P(R)\) contains 0 or a minuscule weight.

